\chapter{Kinetic Theory: Classical, Relativistic, and Quantum}
\section{Math}
\subsection{Eigenfunction expansion and resolvent}
An example of the application of the Liouville equation integrating factor is an ideal gas.

Taking the purely kinetic Hamiltonian

\[
H = \sum_{s=1}^N \frac{p^2_s}{2m}, \quad 0 \leq x^{i}_s \leq L
  \]
  Where $ x^{i}_s$ denotes the  Cartesian component of $\mathbf{x}_s$. Let us call the $\vu{\Lambda}$ operator. Thus:

\[
\vu{\Lambda}_0 = -i \sum_s \pdv{H}{\mathbf{p}_s} \cdot \pdv{}{\mathbf{x}_s} = -i \sum_s \mathbf{v}_s \cdot \pdv{}{\mathbf{x}_s}
  \]

  The Eigenfunctions of $\vu{\Lambda}$ satisfy the eigenvalue equation.

Taking this and solving we get
\[
D(\mathbf{x}^n, \mathbf{p}^N, t) = \frac{1}{L^{3N/2}}\sum_k D_k \mathbf{p}^Ne^{i \sum_s k_s \cdot (\mathbf{x}_s-\mathbf{v}t)}
  \]

  Taking a free particle propagator we find that one evolves the function through time via
\[
  \exp{-t \mathbf{v_s}} \cdot \pdv{}{\mathbf{x}_s} \]

One last solution technique presented to initial value problems for the Liouville equation uses the Laplance transform operated upon the Liouville equation
\[
\int_0^\infty dte^{-st}\left(i \pdv{D}{t} = \vu{\Lambda}D  \right)
  \]

  From there we define
\[
\vu{R} = (\vu{\Lambda} - is)^{-1}
  \]
  as the resolvent operator.

\subsection{Distribution Functions}

Nothing

\subsection{Markov Process}

Given the phase vector $z = (1,2,...,N)$, in terms of this variable, the two-time, or conditional, distribution function may be written at the product
\[
\prod (z,t \mid z_0, t_0) dz
  \]

  This distribution has three fundemental principles
\begin{itemize}
  \item The system undergoes a transition from some state in the interval $t-t_0$
  \item The system does not change its state in zero time
  \item The function $\Pi$ obeys the chapman-kolmogorov equation
\end{itemize}
\[
  \Pi (z,t \mid z_0, t_0) = \int  \Pi (z,t \mid z', t') \Pi (z',t' \mid z_0, t_0)dz'
\]
The integral represents a sum over intermediate states.

For systems whose behavior is homogeneous in time, processes can only depend on the time interval and not the initial or final times.

Our aim is then to obtain a differential equation for $\Pi$ appropriate to the case where phase variables go over a countable set.

The integral equation becomes
\[
\Pi( l \mid l_0 ; t+ \delta t ) = \sum_{\forall j} \Pi ( l \mid j ; \delta t ) \Pi ( j \mid l_0 ; t )
  \]

  Conditions 1 and two become
\[
\sum_{\forall l}\Pi (l \mid l_0; t) =1
  \]
\[
\Pi(l \mid l_0;0) = \delta_{ll_0}
  \]
  For $l \neq j$
\[
\Pi (l \mid j; \delta t) \to 0 \, \text{as} \, \delta t \to 0
  \]

  With some further derivation I won't go into we get the canonical form of the mater equation
\[
  \pdv{\Pi(l \mid l_0 ; t)}{t} = \sum_j (w_{jl} \Pi (j \mid l_0 ; t) - w_{lj} \Pi (l \mid l_0; t))
\]

Which is essentially this psuedo equation
$$\text{Change} = \text{Gain} - \text{Loss}$$

\subsection{Central Limit Theorem}

Variance is writen as
\[
D(\xi) = \sum_{\forall x}(x - \langle x \rangle)^2 P(x)
  \]

  The characteristic function being
\[
  \phi(a) = \sum_{\forall x} P(x)e^{iax}  \]

From here we can take a Fourier transformation expand into a series that will be used later. We can then define a random walk.

An importaint distribution is defined as
\[
P(l,n) = \frac{n !}{r!(n-r)!}\mathfrak{p}^r(1- \mathfrak{p})^{n-r}
  \]
  When we define $r \equiv \frac{n+l}{2}$.

This lets us use some math to define
\[
    P(r) \equiv = P(r, \frac{\lambda}{n} = \frac{\lambda '}{r!}e^{-\lambda}  \]

More commonly known as the Poisson distribution.

\section{Analyses of the Liouville Equation}

\subsection{BBKGY Hierarchy}

Generalizing what I have already done.

\subsection{Correlation Extrapolation: Vlasov limit}
If we assume $N \gets s$, then we can rewrite as
\[
\left(\pdv{}{t} - \vu{L_s}  \right)f_s = \vu{I}_s f_{s+1}
  \]

  With some further math we can define both how dense the fluid is and how strongly coupled.

  Vlasov equation results in the a well defined limit, that can be caused by plasma comprised of electrons in a neutralizing, fixed, positive background.

\subsection{Diagrams: Prigogine Analysis}
